\documentclass[11pt]{article}
\usepackage[a4paper,margin=25mm]{geometry}
\usepackage{amsmath,amssymb,mathtools}
\usepackage{fontspec,unicode-math}
\setmainfont{Latin Modern Roman}
\setsansfont{Latin Modern Sans}
\setmonofont{Latin Modern Mono}
\setmathfont{Latin Modern Math}
\usepackage{xurl,hyperref}
\hypersetup{colorlinks=true,urlcolor=blue,linkcolor=blue}
\newcommand{\sref}[2]{\hyperlink{source:#1:#2}{[S#2]}}
\newcommand{\cataloguescope}[1]{\paragraph{Mathematical question.}#1}
\begin{document}
\section{Globalizing prescribed local Brauer classes}\label{statement}

\subsection{Definitions and mathematical statement}

\paragraph{Shared notation.}$\mathbb N=\{1,2,\ldots\}$, and $\mathbb Z,\mathbb Q,\mathbb R,\mathbb C$ denote the integers, rationals, reals and complex numbers. Any other conventions are specified locally.


A global field $K$ is a number field or a one-variable function field over a finite field. For a place $v$, write $K_v$ for its completion. The Brauer group $\operatorname{Br}(K)$ consists of equivalence classes of finite-dimensional central simple $K$-algebras, with tensor product as addition; $\operatorname{Br}(K)[n]$ denotes its subgroup killed by $n$. Write $\operatorname{inv}_v$ for the local Brauer invariant, with its usual values in $\mathbb Q/\mathbb Z$ (and the real or complex subgroup at archimedean places).
\paragraph{Statement recorded in the supplied source.}
Given $K$, a positive integer $n$, a finite set $S$ of places of $K$, and for each $v\in S$ a class $c_v\in\operatorname{Br}(K_v)[n]$ represented by a central simple algebra $A_v$, explicitly construct one or more classes $c\in\operatorname{Br}(K)$, represented by central simple $K$-algebras $A$, such that
\[
\operatorname{inv}_v(c)=\operatorname{inv}_v(c_v)\quad(v\in S).
\]
Equivalently the localizations $A\otimes_KK_v$ have the prescribed Brauer classes. In particular, address the case in which the local splitting fields are dihedral extensions. \sref{statement}{2}

\subsection{Short English statement}
Explicitly globalize finitely many prescribed local Brauer classes, especially those split by dihedral extensions.
\subsection{Sources}
\begin{description}
\item[\hypertarget{source:statement:1}{S1}] ULAM.AI and the UnsolvedMath Contributors, UnsolvedMath: A Curated Collection of Open Mathematics Problems, catalogue entry. CC BY 4.0. \url{https://huggingface.co/datasets/ulamai/UnsolvedMath}.
\item[\hypertarget{source:statement:2}{S2}] Gerhard Frey, \emph{Duality in Arithmetic Geometry and Applications to Discrete Logarithms}, IPAM Number Theory and Cryptography: Open Problems workshop (9--13 October 2006), Problem 7, PDF pages 45--49. \url{https://helper.ipam.ucla.edu/publications/SCWS1/SCWS1_6219.pdf}.
\end{description}
\label{end:statement}
\end{document}
